% SEGMENT OLP-0074-S01
% Part: first-order-logic
% Chapter: sequent-calculus
% Section: derivations

\documentclass[../../../include/open-logic-section]{subfiles}

\begin{document}

\iftag{FOL}
      {\olfileid{fol}{seq}{der}}
      {\olfileid{pl}{seq}{der}}

\olsection{\usetoken{P}{derivation}}

\begin{explain}
ஒரு தொடக்கத் தொடரணி எப்படி இருக்கும் என்று கூறி, அனுமான விதிகளையும்
வழங்கியுள்ளோம். தொடரணிக் கணியத்தின் !!^{derivation}s இவற்றிலிருந்து
மீள்வரையறையாக உருவாக்கப்படுகின்றன: ஒவ்வொரு !!{derivation}வும் தனித்த ஒரு
தொடக்கத் தொடரணியாக இருக்கும்; அல்லது ஓர் அனுமானத்தைத் தொடர்ந்து வரும் ஒன்று
அல்லது இரண்டு !!{derivation}s ஐக் கொண்டிருக்கும்.
\end{explain}

% SEGMENT OLP-0074-S02
\begin{defn}[$\Log{LK}$ !!{derivation}]
ஒரு தொடரணி~$S$ இன் \emph{$\Log{LK}$-!!{derivation}} என்பது பின்வரும்
நிபந்தனைகளை நிறைவேற்றும் தொடரணிகளின் முடிவுறு மரம்:
\begin{enumerate}
\item மரத்தின் உச்சியிலுள்ள தொடரணிகள் தொடக்கத் தொடரணிகள்.
\item மரத்தின் அடியிலுள்ள தொடரணி~$S$.
\item மரத்தில் $S$ தவிர உள்ள ஒவ்வொரு தொடரணியும், அந்தத் தொடரணிக்கு நேராகக்
  கீழே நிற்கும் முடிவைக் கொண்ட ஓர் அனுமான விதியின் சரியான பயன்பாட்டுக்கு
  முன்கோளாக இருக்கும்.
\end{enumerate}
அப்போது $S$ ஐ அந்த !!{derivation}[gen] \emph{இறுதித் தொடரணி} என்றும், $S$
\emph{$\Log{LK}$ இல் !!{derivable}} (அல்லது $\Log{LK}$-!!{derivable}) என்றும்
கூறுகிறோம்.
\end{defn}

% SEGMENT OLP-0074-S03
\begin{ex}
ஒவ்வொரு தொடக்கத் தொடரணியும், எடுத்துக்காட்டாக $!C \Sequent !C$, ஒரு
!!a{derivation}. இதற்கு, எடுத்துக்காட்டாக, $\LeftR{\Weakening}$ விதியைப்
பயன்படுத்தி ஒரு புதிய !!{derivation}[acc] பெறலாம்:
\begin{prooftree}
\Axiom$ \Gamma \fCenter \Delta $
\RightLabel{\LeftR{\Weakening}}
\UnaryInf$ !A, \Gamma \fCenter \Delta$
\end{prooftree}
ஆனால் இந்த விதி பொதுவானது: விதியிலுள்ள $!A$ ஐ எந்தவொரு !!{sentence}[ins],
எடுத்துக்காட்டாக~$!D$ ஆலும், மாற்றலாம். அதன் முன்கோள் நமது தொடக்கத்
தொடரணி $!C \Sequent !C$ உடன் பொருந்தினால், $\Gamma$, $\Delta$ ஆகிய இரண்டும்
வெறுமனே~$!C$ என்றும், முடிவு $!D, !C \Sequent !C$ என்றும் அமையும். ஆகவே
பின்வருவது !!a{derivation}:
\begin{prooftree}
\Axiom$ !C \fCenter !C $
\RightLabel{\LeftR{\Weakening}}
\UnaryInf$ !D, !C \fCenter !C$
\end{prooftree}
இப்போது இடப்புறத்திலுள்ள இரண்டு !!{sentence}s[acc] இடமாற்ற அனுமதிக்கும்
$\LeftR{\Exchange}$ போன்ற மற்றொரு விதியைப் பயன்படுத்தலாம். ஆகவே பின்வருவதும்
சரியான !!{derivation}:
\begin{prooftree}
\Axiom$ !C \fCenter !C $
\RightLabel{\LeftR{\Weakening}}
\UnaryInf$ !D, !C \fCenter !C$
\RightLabel{\LeftR{\Exchange}}
\UnaryInf$ !C, !D \fCenter !C$
\end{prooftree}
இந்த விதியின் பயன்பாட்டில், விதி
\begin{prooftree}
\Axiom$ \Gamma, !A, !B, \Pi \fCenter \Delta $
\RightLabel{\LeftR{\Exchange}}
\UnaryInf$ \Gamma, !B, !A, \Pi \fCenter \Delta,$
\end{prooftree}
என்று வழங்கப்பட்டது; $\Gamma$, $\Pi$ ஆகிய இரண்டும் வெறுமையாகவும், $\Delta$
என்பது $!C$ ஆகவும், $!A$, $!B$ ஆகியவற்றின் பங்குகளை முறையே $!D$, $!C$
ஆகியவை வகிப்பனவாகவும் இருந்தன. இதேபோன்ற முறையில்,
\begin{prooftree}
\Axiom$ !D \fCenter !D $
\RightLabel{\LeftR{\Weakening}}
\UnaryInf$ !C, !D \fCenter !D$
\end{prooftree}
என்பதும் !!a{derivation} என்று காண்கிறோம்.

% SEGMENT OLP-0074-S04
இப்போது இந்த இரண்டு !!{derivation}s[acc] எடுத்து $\RightR{\land}$ கொண்டு
இணைக்கலாம். அந்த விதி:
\begin{prooftree}
\Axiom$\Gamma \fCenter \Delta, !A$
\Axiom$ \Gamma \fCenter \Delta, !B$
\RightLabel{\RightR{\land}}
\BinaryInf$ \Gamma \fCenter \Delta, !A \land !B $
\end{prooftree}
நமது நேர்வில், முன்கோள்களுடன் முடியும் !!{derivation}s[gen] இறுதித்
தொடரணிகளோடு விதியின் முன்கோள்கள் பொருந்த வேண்டும். இதன் பொருள் $\Gamma$
என்பது $!C, !D$, $\Delta$ வெறுமை, $!A$ என்பது $!C$, $!B$ என்பது $!D$.
ஆகவே அனுமானம் சரியாக இருக்க வேண்டுமெனில், அதன் முடிவு $!C, !D \Sequent !C
\land !D$.
\begin{prooftree}
\Axiom$ !C \fCenter !C $
\RightLabel{\LeftR{\Weakening}}
\UnaryInf$ !D, !C \fCenter !C$
\RightLabel{\LeftR{\Exchange}}
\UnaryInf$ !C, !D \fCenter !C$
\Axiom$ !D \fCenter !D $
\RightLabel{\LeftR{\Weakening}}
\UnaryInf$ !C, !D \fCenter !D$
\RightLabel{\RightR{\land}}
\BinaryInf$ !C, !D \fCenter !C \land !D $
\end{prooftree}
முன்கோள்களை மறுவரிசையிலும் அமைக்கலாம்; அப்போது $!A$ என்பது $!D$ ஆகவும்,
$!B$ என்பது~$!C$ ஆகவும் இருக்கும்.
\begin{prooftree}
\Axiom$ !D \fCenter !D $
\RightLabel{\LeftR{\Weakening}}
\UnaryInf$ !C, !D \fCenter !D$
\Axiom$ !C \fCenter !C $
\RightLabel{\LeftR{\Weakening}}
\UnaryInf$ !D, !C \fCenter !C$
\RightLabel{\LeftR{\Exchange}}
\UnaryInf$ !C, !D \fCenter !C$
\RightLabel{\RightR{\land}}
\BinaryInf$ !C, !D \fCenter !D \land !C $
\end{prooftree}
\end{ex}

\end{document}
